\begin{description}
\item[(a)] The declination is simply the complement of the polar angle:
  \begin{align*}
    \delta &= 90^\circ - \theta\\
    \therefore \delta_\odot &= 21^\circ 30' 15.9''\\
    \delta_{\mathrm{Moon}} &= 21^\circ 12' 18.4''
  \end{align*}
  \hfill(1.0 + 1.0 + 1.0)

\item[(b)] The local sidereal time at Greenwich corresponds by definition to
  the meridian of right ascension that is right above Greenwich. Since
  Greenwich corresponds to an azimuthal angle of $0^h$, the right ascension at
  any given point corresponds simply to the sum of the local sidereal time at
  Greenwich and the azimuthal angle:
  \begin{align*}
    \alpha &= \varphi + LST_{\mathrm{Greenwich}}\\
    \therefore \alpha_\odot &= 4^h 21^m 7.3^s\\
    \alpha_{\mathrm{Moon}} &= 4^h 21^m 12.6^s
  \end{align*}
  \hfill(1.0 + 1.0 + 1.0)

\item[(c)] By subtracting the Cartesian coordinates of the Sun from the
  Cartesian coordinates of the Moon, it is possible to obtain a vector for the
  direction of the axis of the Moon's shadow cone:
  \hfill(1.0)
  \begin{align*}
    \vec{u} &= \langle -1.339 \times 10^{11},\ 4.317 \times 10^{10},\
                       -5.544 \times 10^{10} \rangle\ \mathrm{m}\\
    |\vec{u}| &= 1.512 \times 10^{11}\ \mathrm{m}\\
    \therefore \hat{u} &= \langle -0.8855,\ 0.2855,\ -0.3666 \rangle
  \end{align*}
  \hfill(1.0 + 1.0 + 1.0)

\item[(d)] If one draws a vector in the same direction as $\hat{u}$ up to the
  earth's surface, then its magnitude can be found by the formula, where
  $\vec{M}$ is the position vector of the Moon, $k$ is a constant, and
  $R_\oplus$ is the radius of the Earth:
  \[
    |\vec{M} + k\hat{u}|^2 = R_\oplus^2
  \]
  \hfill(5.0)

  Solving for $k$:
  \begin{align*}
    0 &= k^2 (\hat{u} \cdot \hat{u}) + 2k(\hat{u} \cdot \vec{M})
         + \vec{M} \cdot \vec{M} - R_\oplus^2\\
      &= k^2 + 2k(\hat{u} \cdot \vec{M}) + |\vec{M}|^2 - R_\oplus^2
  \end{align*}
  Now, $|\vec{M}|^2 - R_\oplus^2 = 1.288 \times 10^{17}\ \mathrm{m}$ % sic: units given as m in the source
  \begin{align*}
    2(\hat{u} \cdot \vec{M}) &= -7.178 \times 10^{8}\ \mathrm{m}\\
    \therefore 0 &= k^2 - k(7.178 \times 10^{8}) + (1.288 \times 10^{17})\\
    k &= \frac{7.178 \pm \sqrt{(-7.178)^2 - 4 \times 12.88}}{2} \times 10^{8}
  \end{align*}
  The vector will intersect the earth's surface at two points. Hence two
  solutions. However, only the first intersection with the sphere is relevant
  in this case, so only the smallest solution to the equation should be
  considered:
  \begin{align*}
    k &= \frac{7.178 - \sqrt{(-7.178)^2 - 4 \times 12.88}}{2} \times 10^{8}\\
      &= 3.528 \times 10^{8}
  \end{align*}
  \hfill(4.0)

  The vector of location of the greatest eclipse corresponds to the position
  vector of the Moon plus $k$ times the direction vector $\hat{u}$:
  \begin{align*}
    \vec{p} &= \vec{M} + k\hat{u}\\
      &= \langle 6.091 \times 10^{6},\ -1.828 \times 10^{6},\
                 4.854 \times 10^{5} \rangle
  \end{align*}
  \hfill(2.0)

  It is possible to use the following expressions to convert this vector to
  the spherical system:
  \begin{align*}
    r &= \sqrt{x^2 + y^2 + z^2}\\
      &= 6.378 \times 10^{6}\ \mathrm{m}\\
    \theta &= \arctan\left(\frac{\sqrt{x^2 + y^2}}{z}\right)
      \quad \text{(valid for } z > 0\text{)}\\
      &= 1.495\ \mathrm{rad} = 85^\circ 38'\\
    \varphi &= \arctan\left(\frac{y}{x}\right)
      \quad \text{(valid for } x > 0\text{)}\\
      &= -0.2915\ \mathrm{rad} = -16^\circ 42'
  \end{align*}
  \hfill(2.0)

  The latitude corresponds to the complement of the angle $\theta$, and the
  longitude corresponds to the angle $\varphi$:
  \[
    \boxed{4^\circ 22'\,\mathrm{N},\ 16^\circ 42'\,\mathrm{W}}
  \]
  \hfill(2.0)

\item[(e)] The approximate geometry described in the problem statement is
  shown below:

% FIGURE NEEDED: cross-section of the shadow cone: Sun of radius R_Sun and
% Moon of radius R_Moon with half-angle beta marked as 2*beta at both, the
% umbra of radius 2 r_umbra at distance d_Moon - R_Earth beyond the Moon,
% and the distance d_Sun - d_Moon between Sun and Moon;
% 2024 theory solutions PDF, page 27 (part e).
  \hfill(3.0)

  The angle $\beta$ can be calculated as follows:
  \begin{align*}
    \beta &= \arccos\left(
        \frac{R_\odot - R_{\mathrm{Moon}}}{d_\odot - d_{\mathrm{Moon}}}\right)\\
      &= \arccos\left(
        \frac{6.955 \times 10^{8} - 1.737 \times 10^{6}}
             {1.516 \times 10^{11} - 3.589 \times 10^{8}}\right)\\
      &= 89.74^\circ
  \end{align*}
  \hfill(3.0)

  Using this angle, it is possible to determine the radius of the umbra:
  \[
    r_{\mathrm{umbra}} = \frac{R_{\mathrm{Moon}}
      - (d_{\mathrm{Moon}} - R_\oplus)\cdot\cos(\beta)}{\sin(\beta)}
  \]
  \hfill(3.0)
  \[
    \boxed{r_{\mathrm{umbra}} = 1.196 \times 10^{5}\ \mathrm{m}}
  \]
  \hfill(1.0)

\item[(f)] The rotational velocity of the Earth at the latitude of the center
  of the umbra is the following:
  \begin{align*}
    v_{\mathrm{rot}} &= \frac{2\pi R_\oplus}{23^h 56^m 04^s}
        \cos(\phi_{\mathrm{umbra}})\\
      &= \frac{2\pi \times 6.378 \times 10^{6}}{23^h 56^m 04^s}
        \cos(4^\circ 22')
  \end{align*}
  \hfill(2.0)
  \[
    \boxed{v_{\mathrm{rot}} = 463.7\ \mathrm{m/s}}
  \]
  \hfill(1.0)

\item[(g)] Using the vis-viva equation:
  \begin{align*}
    v_{\mathrm{Moon}} &= \sqrt{GM_\oplus\left(\frac{2}{d_{\mathrm{Moon}}}
        - \frac{1}{a_{\mathrm{Moon}}}\right)}\\
      &= \sqrt{6.674 \times 10^{-11} \times 5.972 \times 10^{24}
        \left(\frac{2}{3.589 \times 10^{8}}
            - \frac{1}{3.844 \times 10^{8}}\right)}
  \end{align*}
  \hfill(2.0)
  \[
    \boxed{v_{\mathrm{Moon}} = 1088\ \mathrm{m/s}}
  \]
  \hfill(2.0)

\item[(h)] It is possible to use the following diagram to determine the value
  of the angle $\kappa$ between the $x$-axis and the velocity vector of the
  Moon:

% FIGURE NEEDED: spherical triangle on the celestial sphere with the Moon at
% the top vertex, the celestial equator along the bottom through the vernal
% point (gamma), the lunar orbit crossing the equator, sides delta_Moon and
% alpha_Moon marked, the arc 24h - alpha_intersection, right angles at the
% Moon's meridian, and the angle kappa at the Moon;
% 2024 theory solutions PDF, page 29 (part h).
  \hfill(6.0)

  Note that the great circle that intersects the Moon's meridian of right
  ascension at a right angle has a tangent line that coincides with the
  $x$-axis, so the angle $\kappa$ on the figure is also the angle between the
  velocity vector of the Moon and the $x$-axis. Also note that the Moon is
  moving from west to east in the celestial sphere (to the left in the
  figure), so its velocity vector is on the first quadrant of the $xy$-plane
  on this system.

  Using the four parts formula:
  \begin{align*}
    \cos(\delta_{\mathrm{Moon}})\cos(90^\circ)
      &= \sin(\delta_{\mathrm{Moon}})
         \cot(\alpha_{\mathrm{Moon}} + 24^h - \alpha_{\mathrm{intersection}})\\
      &\quad - \sin(90^\circ)\cot(90^\circ - \kappa)
  \end{align*}
  \hfill(3.0)
  \begin{align*}
    \therefore 0 &= \sin(\delta_{\mathrm{Moon}})
        \cot(\alpha_{\mathrm{Moon}} - \alpha_{\mathrm{intersection}})
        - \tan(\kappa)\\
    \tan(\kappa) &= \frac{\sin(\delta_{\mathrm{Moon}})}
        {\tan(\alpha_{\mathrm{Moon}} - \alpha_{\mathrm{intersection}})}\\
    \kappa &= \arctan\left[\frac{\sin(\delta_{\mathrm{Moon}})}
        {\tan(\alpha_{\mathrm{Moon}} - \alpha_{\mathrm{intersection}})}\right]\\
      &= \arctan\left[\frac{\sin(21^\circ 12' 18.4'')}
        {\tan(4^h 21^m 12.6^s - 23^h 07^m 59.2^s)}\right]\\
    \kappa &= 4^\circ 17'
  \end{align*}
  \hfill(2.0)

  Now it is possible to break down the velocity of the Moon into the $x$ and
  $y$ components:
  \[
    \mathbf{v}_{\mathrm{II}} =
    \begin{bmatrix}
      v_{\mathrm{Moon}}\cos(\kappa)\\
      v_{\mathrm{Moon}}\sin(\kappa)\\
      0
    \end{bmatrix}
    =
    \begin{bmatrix}
      1085\ \mathrm{m/s}\\
      81.25\ \mathrm{m/s}\\
      0
    \end{bmatrix}
  \]
  \hfill(3.0)

  The radial component of the velocity was neglected, so the velocity on the
  $z$-axis is equal to zero.

\item[(i)] It is possible to obtain system III by rotating system II from
  north to south by an angle of
  $\theta_{\mathrm{umbra}} - \theta_{\mathrm{Moon}}$ around the $x$-axis. The
  angles $\theta_{\mathrm{umbra}}$ and $\theta_{\mathrm{Moon}}$ correspond to
  the polar angles of the umbra and the Moon on system I.

  The following two figures illustrate this rotation.

% FIGURE NEEDED: two diagrams: (left) a circle (the Earth) with the System II
% axes drawn at the Moon's polar angle theta_Moon and the System III axes at
% theta_umbra; (right) the axes of Systems II and III sharing the x-axis, with
% the y and z axes rotated by theta_umbra - theta_Moon;
% 2024 theory solutions PDF, page 30 (part i).
  \hfill(3.0)

  Since the rotation is around the $x$-axis, the $x$ component remains
  unchanged. The component that was originally in the $y$-axis is broken down
  into two components on the $y$-axis and the $z$-axis of the new system,
  which results in the following vector:
  \[
    \mathbf{v}_{\mathrm{III}} =
    \begin{bmatrix}
      v_{\mathrm{Moon}} \cdot \cos(\kappa)\\
      v_{\mathrm{Moon}} \cdot \sin(\kappa) \cdot
        \cos(\theta_{\mathrm{umbra}} - \theta_{\mathrm{Moon}})\\
      -\,v_{\mathrm{Moon}} \cdot \sin(\kappa) \cdot
        \sin(\theta_{\mathrm{umbra}} - \theta_{\mathrm{Moon}})
    \end{bmatrix}
  \]
  \hfill(6.0)
  \[
    \mathbf{v}_{\mathrm{III}} =
    \begin{bmatrix}
      1085\ \mathrm{m/s}\\
      77.76\ \mathrm{m/s}\\
      -23.54\ \mathrm{m/s}
    \end{bmatrix}
  \]
  \hfill(1.0)

\item[(j)] The velocity of the umbra corresponds to the vector sum of the
  component caused by the rotation of the Earth and the component caused by
  the velocity of the Moon.

  Note that all points on the axis of the Moon's shadow cone have the same
  velocity, $\mathbf{v}_{\mathrm{III}}$ is also the velocity of the axis of
  the Moon's shadow cone at the center of the umbra. Since the umbra moves
  along the surface of the Earth, so the $z_{\mathrm{III}}$-component should
  be set to zero.
  \[
    \mathbf{v}_{\mathrm{umbra}} =
    \begin{bmatrix}
      v_{\mathrm{III},x} - v_{\mathrm{rot}}\\
      v_{\mathrm{III},y}\\
      0
    \end{bmatrix}
    =
    \begin{bmatrix}
      1085 - 463.7\\
      77.76\\
      0
    \end{bmatrix}
    =
    \begin{bmatrix}
      621.4\ \mathrm{m/s}\\
      77.59\ \mathrm{m/s}\\ % sic: 77.59 here, though the middle step and the modulus below use 77.76
      0
    \end{bmatrix}
  \]
  \hfill(4.0)

  The modulus of this vector is the following:
  \[
    v_{\mathrm{umbra}} = |\mathbf{v}_{\mathrm{umbra}}|
      = \sqrt{621.4^2 + 77.76^2} = 626.3\ \mathrm{m/s}
  \]
  \hfill(2.0)

\item[(k)] Considering the assumptions listed at the beginning of part II, it
  is possible to estimate the length of the totality by dividing the diameter
  of the umbra by its speed:
  \[
    \Delta t_{\mathrm{totality}}
      = \frac{2 \cdot r_{\mathrm{umbra}}}{v_{\mathrm{umbra}}}
      = \frac{2 \times 1.196 \times 10^{5}}{626.3}
  \]
  \hfill(2.0)
  \[
    \boxed{\Delta t_{\mathrm{totality}} = 6\,\mathrm{min}\,21.4\,\mathrm{s}}
  \]
  \hfill(1.0)
\end{description}
